\documentclass[10pt,a4paper]{amsart}
\usepackage{amsmath,amssymb,amsthm,graphicx}
% Retained minimal preamble for the Stacks Project chapter topology.tex (Topology).
% The Stacks Project uses one single shared preamble for every chapter; the snapshot stages it
% identically as sources/stacks-topology/preamble.tex, sources/stacks-chow/preamble.tex and
% sources/stacks-etale-cohomology/preamble.tex.  The block below is copied mechanically, line for
% line, from that shared preamble: the Xy-pic setup (lines 14-18) and the native theorem
% declarations together with the mathematical macros (lines 154-217), plus the project's own
% reference/slogan/history display environments (lines 8-12, via the verbatim package).
% Nothing was added, removed or redefined; the wrapper class and the amsmath/amssymb/amsthm
% base packages are supplied by the assembler.
% Chapter: topology.tex -- Topology.
\usepackage[all]{xy}

% We use 2cell for 2-commutative diagrams.
\xyoption{2cell}

% Theorem environments.
%
\theoremstyle{plain}
\newtheorem{theorem}[subsection]{Theorem}
\newtheorem{proposition}[subsection]{Proposition}
\newtheorem{lemma}[subsection]{Lemma}

\theoremstyle{definition}
\newtheorem{definition}[subsection]{Definition}
\newtheorem{example}[subsection]{Example}
\newtheorem{exercise}[subsection]{Exercise}
\newtheorem{situation}[subsection]{Situation}

\theoremstyle{remark}
\newtheorem{remark}[subsection]{Remark}
\newtheorem{remarks}[subsection]{Remarks}

\numberwithin{equation}{subsection}

% Macros
%
\def\lim{\mathop{\mathrm{lim}}\nolimits}
\def\colim{\mathop{\mathrm{colim}}\nolimits}
\def\Spec{\mathop{\mathrm{Spec}}}
\def\Hom{\mathop{\mathrm{Hom}}\nolimits}
\def\Ext{\mathop{\mathrm{Ext}}\nolimits}
\def\SheafHom{\mathop{\mathcal{H}\!\mathit{om}}\nolimits}
\def\SheafExt{\mathop{\mathcal{E}\!\mathit{xt}}\nolimits}
\def\Sch{\mathit{Sch}}
\def\Mor{\mathop{\mathrm{Mor}}\nolimits}
\def\Ob{\mathop{\mathrm{Ob}}\nolimits}
\def\Sh{\mathop{\mathit{Sh}}\nolimits}
\def\NL{\mathop{N\!L}\nolimits}
\def\CH{\mathop{\mathrm{CH}}\nolimits}
\def\proetale{{pro\text{-}\acute{e}tale}}
\def\etale{{\acute{e}tale}}
\def\QCoh{\mathit{QCoh}}
\def\Ker{\mathop{\mathrm{Ker}}}
\def\Im{\mathop{\mathrm{Im}}}
\def\Coker{\mathop{\mathrm{Coker}}}
\def\Coim{\mathop{\mathrm{Coim}}}

% Boxtimes
%
\DeclareMathSymbol{\boxtimes}{\mathbin}{AMSa}{"02}

%
% Macros for moduli stacks/spaces
%
\def\QCohstack{\mathcal{QC}\!\mathit{oh}}
\def\Cohstack{\mathcal{C}\!\mathit{oh}}
\def\Spacesstack{\mathcal{S}\!\mathit{paces}}
\def\Quotfunctor{\mathrm{Quot}}
\def\Hilbfunctor{\mathrm{Hilb}}
\def\Curvesstack{\mathcal{C}\!\mathit{urves}}
\def\Polarizedstack{\mathcal{P}\!\mathit{olarized}}
\def\Complexesstack{\mathcal{C}\!\mathit{omplexes}}
% \Pic is the operator that assigns to X its picard group, usage \Pic(X)
% \Picardstack_{X/B} denotes the Picard stack of X over B
% \Picardfunctor_{X/B} denotes the Picard functor of X over B
\def\Pic{\mathop{\mathrm{Pic}}\nolimits}
\def\Picardstack{\mathcal{P}\!\mathit{ic}}
\def\Picardfunctor{\mathrm{Pic}}
\def\Deformationcategory{\mathcal{D}\!\mathit{ef}}

% --- appended 2026-10-12 (Chief): the project own definitions of the reference/slogan/history
% display environments, copied mechanically from the shared Stacks preamble lines 8-12
% (\usepackage{verbatim} plus the three \newenvironment lines, which the chapter preamble keeps
% verbatim from the project preamble). Nothing added, removed or redefined: these environments
% carry no mathematical content and the project itself typesets them as comments. Without these
% three lines a counted statement whose own text uses them fails with Environment slogan undefined.
\usepackage{verbatim}
\newenvironment{reference}{\comment}{\endcomment}
\newenvironment{slogan}{\comment}{\endcomment}
\newenvironment{history}{\comment}{\endcomment}
\title{Alexander's subbase theorem: a finite-refinement test on a subbase implies quasi-compactness}
\author{The Stacks Project Authors}
\date{stacks-project:topology.tex:local-snapshot-20261012; CC BY 4.0}
\begin{document}
\maketitle
\noindent\textit{Editorial note.} All delivered mathematical text is the author's own, copied from the pinned Stacks Project chapter source sources/stacks-topology/pinned-original-source.tex; no statement or proof step was rewritten, reordered or removed.\par\medskip
\noindent\textit{Editorial note.} Preamble provenance: the Stacks Project keeps one single shared preamble for every chapter, staged here as sources/stacks-topology/preamble.tex; this card uses retained-preamble.tex, which copies that shared preamble mechanically (the Xy-pic setup, the native theorem declarations, the mathematical macros and the reference/slogan/history display environments). Nothing was added, removed or redefined; the wrapper class and the base packages are supplied by the assembler.\par\medskip
\noindent\textit{Editorial note.} Bibliography provenance: the chapter carries no chapter-local bibliography (the project keeps one shared bibliography file that is not staged here) and neither the counted statement nor its proof contains any citation, so the card carries no bibliography and no entry was reconstructed.\par\medskip
\noindent\textit{Editorial note.} This card was selected from the staged candidate list for this chapter (candidates-halfA.json / candidates-halfB.json), where it is the candidate labelled lemma-subbase-theorem; the statement and proof spans are the candidate's own statement\_lines and proof\_lines, checked against the source so that each span opens and closes a real, uncommented environment.\par\medskip
\noindent\textit{Allowed standard prerequisites (editorial).}\begin{itemize}
\item The notion of a subbase for a topological space, and the definition of quasi-compactness as the finite-refinement property of open coverings. Set up in the author's sections on quasi-compact spaces and maps, immediately before the counted statement. Fixes the hypothesis on the subbase and the conclusion to be proved.
\end{itemize}\medskip
\begin{lemma}[Alexander subbase theorem]
\label{lemma-subbase-theorem}
Let $X$ be a topological space. Let $\mathcal{B}$ be a subbase for $X$.
If every covering of $X$ by elements of $\mathcal{B}$ has a finite
refinement, then $X$ is quasi-compact.
\end{lemma}

\begin{proof}
Assume there is an open covering of $X$ which does not have a finite
refinement. Using Zorn's lemma we can choose a maximal open covering
$X = \bigcup_{i \in I} U_i$ which does not have a finite refinement
(details omitted).
In other words, if $U \subset X$ is any open which does not occur as
one of the $U_i$, then the covering $X = U \cup \bigcup_{i \in I} U_i$
does have a finite refinement. Let $I' \subset I$ be the set of indices
such that $U_i \in \mathcal{B}$. Then $\bigcup_{i \in I'} U_i \not = X$,
since otherwise we would get a finite refinement covering $X$ by our
assumption on $\mathcal{B}$. Pick $x \in X$,
$x \not \in \bigcup_{i \in I'} U_i$. Pick $i \in I$ with $x \in U_i$.
Pick $V_1, \ldots, V_n \in \mathcal{B}$ such that
$x \in V_1 \cap \ldots \cap V_n \subset U_i$. This is
possible as $\mathcal{B}$ is a subbasis for $X$. Note that
$V_j$ does not occur as a $U_i$. By maximality of the chosen
covering we see that for each $j$ there exist
$i_{j, 1}, \ldots, i_{j, n_j} \in I$ such that
$X = V_j \cup U_{i_{j, 1}} \cup \ldots \cup U_{i_{j, n_j}}$.
Since $V_1 \cap \ldots \cap V_n \subset U_i$ we conclude that
$X = U_i \cup \bigcup U_{i_{j, l}}$ a contradiction.
\end{proof}

\end{document}
